\SourceLeaf{352}
discerptiones propriae formae \(\varphi\) indefinite exhibeantur per
\[
(gt+hu)^2+(g't+h'u)^2+(g''t+h''u)^2
\]
omnes discerptiones formae \(\varphi'\) expressum iri per
\[
(gt-hu)^2+(g't-h'u)^2+(g''t-h''u)^2
\]
nec non ex his easdem discerptiones numeri \(M\) derivari ut ex illis. Denique pro eo casu ubi \(\varphi\) est forma e classe ancipite, attamen neque e classe principali neque formae \((2,0,\frac12 M)\) aut \((2,1,\frac12(M+1))\) aequivalens (prout \(M\) par aut impar), e valoribus expr. \(\surd-(p,-q,r)\) semissem omittere licet; sed brevitatis caussa hocce compendium fusius hic non explicamus. --- Ceterum iisdem compendiis etiam uti possumus, quando omnes repraesentationes propriae ipsius \(M\) per \(xx+yy+zz\) desiderantur, quum hae e discerptionibus facillime evolvantur.

Exempli caussa investigabimus omnes discerptiones numeri \(770\) in terna quadrata, ubi \(\mu=3\), \(e=e'=0\), adeoque \(E=2k\). Per classificationem formarum binariarum positivarum determinantis \(-770\), quam quoniam a quovis ad normam art. 231 facile condi potest brevitatis gratia non adscribimus, invenitur classium positivarum multitudo \(=32\), quae omnes sunt proprie primitivae et inter \(8\) genera distribuuntur, ita ut sit \(k=4\), et proin \(E=8\). Genus, cuius numerus characteristicus \(-1\), respectu numerorum \(5,7,11\) manifesto characteres particulares \(R5;\ N7;\ N11\) habere debet, unde per art. 263 facile concluditur, ipsius characterem respectu numeri \(8\) esse debere \(1\text{ et }3,\ 8\). Iam in eo genere, cuius character \(1\text{ et }3,\ 8;\ R5;\ N7;\ N11\), quatuor classes reperiuntur, pro quarum repraesentantibus eligimus formas \((6,2,129)\), \((6,-2,129)\), \((19,3,41)\), \(19,-3,41)\); classem secundam vero et quartam reiicimus, utpote primae et tertiae oppositas. Quatuor discerptiones formae \((19,3,41)\) iam in art. 289 tradidimus, e quibus sequuntur discerptiones numeri \(770\) in \(9+361+400\), \(16+25+729\), \(81+400+289\), \(576+169+25\). Simili ratione inveniuntur quatuor discerptiones formae \(6tt+4tu+129uu\) in
\[
\begin{gathered}
(t-8u)^2+(2t+u)^2+(t+8u)^2,\\
(t-10u)^2+(2t+5u)^2+(t+2u)^2\\
(2t-5u)^2+(t+10u)^2+(t+2u)^2,\\
(2t+7u)^2+(t-8u)^2+(t-4u)^2
\end{gathered}
\]
\SourceLeaf{353}
resp. e valoribus expressionis \(\surd-(6,-2,129)\) hisce oriundae \((48,369)\), \((62,-149)\), \((92,-159)\), \((202,61)\); unde prodeunt discerptiones numeri \(770\) in \(225+256+289\), \(1+144+625\), \(64+81+625\), \(16+225+529\). Praeter has octo discerptiones aliae non dantur.

Quae ad discerptiones numerorum in terna quadrata divisores communes habentia attinent, tam facile e theoria generali art. 281 sequuntur, ut non opus sit huic rei immorari.

\Needspace{8\baselineskip}
\begin{center}\textit{Demonstratio theorematum Fermatianorum. quemvis integrum in tres numeros trigonales vel quatuor quadrata discerpi posse.}\end{center}
\Article{293}
Disquisitiones praecedentes etiam suppeditant demonstrationem theorematis famosi, \textit{omnem numerum integrum positivum in tres numeros trigonales discerpi posse}, quod a Fermatio olim inventum est, sed cuius demonstratio rigorosa hactenus desiderabatur, Manifestum est, quamvis discerptionem numeri \(M\) in trigonales
\[
\tfrac12 x(x+1)+\tfrac12 y(y+1)+\tfrac12 z(z+1)
\]
producere discerptionem numeri \(8M+3\) in terna quadrata imparia
\[
(2x+1)^2+(2y+1)^2+(2z+1)^2
\]
et vice versa. Quivis autem numerus integer positivus \(8M+3\) per theoriam praecedentem in tria quadrata resolubilis est, quae necessaria erunt imparia (V. annot. art. 291); resolutionumque multitudo pendet tum a multitudine factorum primorum ipsius \(8M+3\), tum a multitudine classium, in quas formae binariae determinantis \(-(8M+3)\) distribuuntur. Totidem discerptiones numeri \(M\) in ternos trigonales dabuntur. Supponimus autem, \(\frac12 x(x+1)\) pro valore quocunque integro ipsius \(x\) tamquam trigonalem spectari; quodsi magis placeret cifram excludere, theorema ita immutare oporteret: Quivis integer positivus vel ipse trigonalis est, vel in duos vel in tres trigonales resolubilis. Similis mutatio in theoremate sequente facienda esset, si cifram a quadratis excludere placeret.

Ex iisdem principiis demonstratur aliud Fermatii theorema, \textit{quemvis numerum integrum positivum in quatuor quadrata decomponi posse.} Subtrahendo a numero formae \(4n+2\) quadratum arbitrarium (illo numero minus), a numero formae
\SourceLeaf{354}
\(4n+1\) quadratum par, a numero formae \(4n+3\) quadratum impar, residuum in omnibus his casibus in tria quadrata resolubile erit, adeoque numerus propositus in quatuor. Denique numerus formae \(4n\) exhiberi potest per \(4^{\mu}N\) ita ut \(N\) ad aliquam trium formarum praecedentium pertineat: resoluto autem ipso \(N\) in quatuor quadrata, etiam \(4^{\mu}N\) resolutus erit. A numero formae \(8n+3\) etiam subduci potest quadratum radicis pariter paris, a numero formae \(8n+7\) quadratum radicis impariter paris, a numero formae \(8n+4\) quadratum impar, residuumque in tria quadrata resolubile erit. Ceterum hocce theorema iam ab ill. La Grange demonstratum erat, \textit{Nouv. Mém. de l'Ac. de Berlin} 1770 \textit{p.} 123, quam demonstrationem (a nostra prorsus diversam) fusius explicavit ill. Euler in \textit{Actis Ac. Petr. Vol.} II. \textit{p.} 48. --- Alia Fermatii theoremata quae praecedentium quasi continuationem constituunt, quemvis numerum integrum in quinque numeros pentagonales, sex hexagonales, septem heptagonales \textit{etc.} resolubilem esse, demonstratione hactenus carent, altiaque principia requirere videntur.

\Needspace{9\baselineskip}
\begin{center}\textit{Solutio aequationis \(axx+byy+czz=0\).}\end{center}
\Article{294}
\textsc{Theorema.} \textit{Designantibus \(a,b,c\) numeros inter se primos, quorum nullus neque \(=0\) neque per quadratum divisibilis, aequatio}
\[
axx+byy+czz=0\ldots(\Omega)
\]
\textit{resolutionem in integris non admittet (praeter hanc \(x=y=z=0\) ad quam non respicimus) nisi \(-bc,-ac,-ab\) resp. sint residua quadratica ipsorum \(a,b,c\), atque hi numeri signis inaequalibus affecti; his vero quatuor conditionibus locum habentibus, \((\Omega)\) in integris resolubilis erit.}

\textit{Dem.} Si \((\Omega)\) per integros omnino est resolubilis, etiam per tales valores ipsorum \(x,y,z\) resolvi poterit, qui divisorem communem non habent; nam valores quicunque aequ. \(\Omega\) satisfacientes, etiamnum satisfacient, si per divisorem communem maximum dividuntur. Iam supponendo \(app+bqq+crr=0\), atque \(p,q,r\) a divisore communi liberos, etiam inter se primi erunt; si enim \(q,r\) divisorem communem \(\mu\) haberent, hic ad \(p\) primus esset, \(\mu\mu\) autem metiretur ipsum \(app\) adeoque etiam ipsum \(a\), contra hyp.; et perinde \(p,r;\ p,q\) inter se primi erunt. Repraesentatur itaque \(-app\) per formam binariam \(byy+czz\),
\SourceLeaf{355}
tribuendo ipsis \(y,z\) valores inter se primos \(q,r\); unde illius determinans \(-bc\) residuum quadraticum ipsius \(app\) adeoque etiam ipsius \(a\) erit (art. 154); eodem modo erit \(-acRb,-abRc\). Quod vero \((\Omega)\) resolutionem admittere non possit, si \(a,b,c\) idem signum habeant, tam obvium est, ut explicatione non egeat.

Demonstrationem propositionis inversae. quae theorematis partem secundam constituit, ita adornabimus, ut \textit{primo} formam ternariam ipsi \(\left(\begin{smallmatrix}a,b,c\\0,0,0\end{smallmatrix}\right)\ldots f\) aequivalentem invenire doceamus, cuius coëfficientes \(2,3,4\) per \(abc\) divisibiles sint, unde \textit{secundo} solutionem aequationis \((\Omega)\) deducemus.

I. Investigentur tres integri \(A,B,C\) a divisore communi liberi, atque ita comparati, ut \(A\) primus sit ad \(b\) et \(c\); \(B\) ad \(a\) et \(c\); \(C\) ad \(a\) et \(b\); \(aAA+bBB+cCC\) autem per \(abc\) divisibilis, quod efficietur sequenti modo. Sint \(\mathfrak A,\mathfrak B,\mathfrak C\) resp. valores expressionum \(\surd-bc\pmod a\), \(\surd-ac\pmod b\), \(\surd-ab\pmod c\), qui necessario ad \(a,b,c\) resp. primi erunt. Accipiantur tres integri \(\mathfrak a,\mathfrak b,\mathfrak c\) omnino ad lubitum, modo ita ut ad \(a,b,c\) resp. primi sint (\textit{e. g.} omnes \(=1\)), determinenturque \(A,B,C\) ita ut sit
\[
\begin{aligned}
A&\equiv\mathfrak b c\pmod b&&\text{et }\equiv\mathfrak c\mathfrak C\pmod c\\
B&\equiv\mathfrak c a\pmod c&&\text{et }\equiv\mathfrak a\mathfrak A\pmod a\\
C&\equiv\mathfrak a b\pmod a&&\text{et }\equiv\mathfrak b\mathfrak B\pmod b
\end{aligned}
\]
Tunc fiet
\[
\begin{aligned}
aAA+bBB+cCC&\equiv\mathfrak a\mathfrak a(b\mathfrak A\mathfrak A+cbb)\\
&\equiv\mathfrak a\mathfrak a(b\mathfrak A\mathfrak A-\mathfrak A\mathfrak A b)\equiv0\pmod a
\end{aligned}
\]
sive per \(a\) divisibilis, et perinde per \(b,c\), adeoque etiam per \(abc\) divisibilis erit. Praeterea patet, \(A\) necessario fieri primum ad \(b\) et \(c\); \(B\) ad \(a\) et \(c\); \(C\) ad \(a\) et \(b\). Si vero hi valores ipsorum \(A,B,C\) divisorem communem (maximum) \(\mu\) implicant, hic manifesto ad \(a,b,c\) adeoque ad \(abc\) primus erit; quare illos valores per \(\mu\) dividendo novos obtinebimus, qui divisorem communem non habebunt, valorem ipsius \(aAA+bBB+cCC\) etiamnum per \(abc\) divisibilem producent, adeoque omnibus conditionibus satisfacient.

II. Numeris \(A,B,C\), hoc modo determinatis. etiam \(Aa,Bb,Cc\) divisorem communem non habebunt. Si enim haberent div. comm. \(\mu\), hic necessario primus esset ad \(a\) (quippe qui tum ad \(Bb\) tum ad \(Cc\) primus est) et similiter ad
\SourceLeaf{356}
\(b\) et \(c\); quare \(\mu\) etiam ipsos \(A,B,C\) metiri deberet, contra hyp. Inveniri poterunt itaque integri \(\alpha,\beta,\gamma\) tales, ut sit \(\alpha Aa+\beta Bb+\gamma Cc=1\); quaerantur insuper sex integri \(\alpha',\beta',\gamma',\alpha'',\beta'',\gamma''\) tales, ut sit
\[
\beta'\gamma''-\gamma'\beta''=Aa,\quad
\gamma'\alpha''-\alpha'\gamma''=Bb,\quad
\alpha'\beta''-\beta'\alpha''=Cc
\]
Iam transeat \(f\) per substitutionem.
\[
\begin{array}{ccc}
\alpha,&\alpha',&\alpha''\\
\beta,&\beta',&\beta''\\
\gamma,&\gamma',&\gamma''
\end{array}
\]
in \(\left(\begin{smallmatrix}m,m',m''\\n,n',n''\end{smallmatrix}\right)=g\) (quae ipsi \(f\) aequivalens erit), dicoque \(m',m'',n\) per \(abc\) divisibiles fore. Ponatur enim
\[
\begin{aligned}
\beta''\gamma-\gamma''\beta&=A',&\gamma''\alpha-\alpha''\gamma&=B',&\alpha''\beta-\beta''\alpha&=C'\\
\beta\gamma'-\gamma\beta'&=A'',&\gamma\alpha'-\alpha\gamma'&=B'',&\alpha\beta'-\beta\alpha'&=C''
\end{aligned}
\]
eritque
\[
\begin{aligned}
\alpha'&=B''Cc-C''Bb,&\beta'&=C''Aa-A''Cc,&\gamma'&=A''Bb-B''Aa\\
\alpha''&=C'Bb-B'Cc,&\beta''&=A'Cc-C'Aa,&\gamma''&=B'Aa-A'Bb
\end{aligned}
\]
Quibus valoribus in aequationibus
\[
\begin{aligned}
m'&=a\alpha'\alpha'+b\beta'\beta'+c\gamma'\gamma'\\
m''&=a\alpha''\alpha''+b\beta''\beta''+c\gamma''\gamma''\\
n&=a\alpha'\alpha''+b\beta'\beta''+c\gamma'\gamma''
\end{aligned}
\]
substitutis, fit, secundum modulum \(a\),
\[
\begin{aligned}
m'&\equiv bcA''A''(BBb+CCc)\equiv0\\
m''&\equiv bcA'A'(BBb+CCc)\equiv0\\
n&\equiv bcA'A''(BBb+CCc)\equiv0
\end{aligned}
\]
i. e. \(m',m'',n\) per \(a\) divisibiles erunt; similique modo iidem numeri per \(b,c\) adeoque etiam per \(abc\) divisibiles inveniuntur. Q. E. P.
\SourceLeaf{357}
\par
III. Ponamus, concinnitatis caussa, determinantem formarum \(f,g\). i. e. numerum \(-abc=d\),
\[
md=M,\quad m'=M'd,\quad m''=M''d,\quad n=Nd,\quad n'=N',\quad n''=N''
\]
patetque, \(f\) transire per substitutionem \((S)\)
\[
\begin{array}{ccc}
\alpha d,&\alpha',&\alpha''\\
\beta d,&\beta',&\beta''\\
\gamma d,&\gamma',&\gamma''
\end{array}
\]
in formam ternariam \(\left(\begin{smallmatrix}Md,M'd,M''d\\Nd,N'd,N''d\end{smallmatrix}\right)=g'\) determinantis \(d^3\), quae itaque sub \(f\) contenta erit. Iam dico, huic formae \(g'\) necessario aequivalere hanc \(\left(\begin{smallmatrix}d,0,0\\d,0,0\end{smallmatrix}\right)=g''\). Patet enim, \(\left(\begin{smallmatrix}M,M',M''\\N,N',N''\end{smallmatrix}\right)=g'''\) fore formam ternariam determinantis \(1\); porro quum per hyp. \(a,b,c\) eadem signa non habeant, \(f\) erit forma indefinita, unde facile concluditur, etiam \(g'\) et \(g'''\) indefinitas esse debere; quare \(g'''\) aequivalebit formae \(\left(\begin{smallmatrix}1,0,0\\1,0,0\end{smallmatrix}\right)\) (art. 277), poteritque transformatio \((S')\) illius in hanc inveniri; manifesto autem per \((S')\) forma \(g'\) transibit in \(g''\). Hinc etiam \(g''\) sub \(f\) contenta erit, et ex combinatione substitutionum \((S),(S')\) deducetur transformatio formae \(f\) in \(g''\). Quae si fuerit
\[
\begin{array}{ccc}
\delta,&\delta',&\delta''\\
\varepsilon,&\varepsilon',&\varepsilon''\\
\zeta,&\zeta',&\zeta''
\end{array}
\]
manifestum est, duplicem solutionem aequationis \((\Omega)\) haberi, puta \(x=\delta',y=\varepsilon',z=\zeta'\), et \(x=\delta'',y=\varepsilon'',z=\zeta''\); simul patet, neutros valores simul \(=0\) evadere posse, quum necessario fiat
\[
\delta\varepsilon'\zeta''+\delta'\varepsilon''\zeta+\delta''\varepsilon\zeta'-\delta\varepsilon''\zeta'-\delta'\varepsilon\zeta''-\delta''\varepsilon'\zeta=d.\quad\text{Q. E. S.}
\]

\emph{Exemplum.} Sit aequatio proposita \(7xx-15yy+23zz=0\), quae resolubilis est, quia \(345\,R\,7,-161\,R\,15,105\,R\,23\). Habentur hic valores ipsorum \(\mathfrak A,\mathfrak B,\mathfrak C\) hi \(3,7,6\); faciendoque \(\mathfrak a=\mathfrak b=\mathfrak c=1\) invenitur \(A=98,B=-39,C=-8\). Hinc eruitur substitutio
\[
\left\{\begin{array}{rrr}3,&5,&22\\-1,&2,&-28\\8,&25,&-7\end{array}\right\}
\]
per quam \(f\) transit in \(\left(\begin{smallmatrix}1520,14490,-7245\\-2415,-1246,4735\end{smallmatrix}\right)=g\). Hinc fit
\[
(S)=\left\{\begin{array}{rrr}7245,&5,&22\\-2415,&2,&-28\\19320,&25,&-7\end{array}\right\},\quad
g'''=\left(\begin{smallmatrix}3670800,6,-3\\-1,-1246,4735\end{smallmatrix}\right).
\]
\SourceLeaf{358}
Forma \(g'''\) transire invenitur in \(\left(\begin{smallmatrix}1,0,0\\1,0,0\end{smallmatrix}\right)\) per substitutionem
\[
\left\{\begin{array}{rrr}3,&5,&1\\-2440,&-4066,&-813\\-433,&-722,&-144\end{array}\right\}\ldots(S')
\]
qua cum \((S)\) combinata prodit haec: \(\left\{\begin{smallmatrix}9,11.12\\-1,9,-9\\-9,4,3\end{smallmatrix}\right\}\) per quam \(f\) transit in \(g''\). Habemus itaque duplicem aequationis propositae solutionem \(x=11,y=9,z=4\), et \(x=12,y=-9,z=3\); posterior simplicior redditur dividendo valores per divisorem communem \(3\), unde \(x=4,y=-3,z=1\).

\Article{295}
Pars posterior theorematis art. praec. etiam sequenti modo absolvi potest. Quaeratur integer \(h\) talis, ut sit \(ah\equiv\mathfrak C\pmod c\), (characteres \(\mathfrak A,\mathfrak B,\mathfrak C\) eadem significatione accipimus ut in art. praec.), fiatque \(ahh+b=ci\). Tunc facile perspicitur, \(i\) fieri integrum, numerumque \(-ab\) esse determinantem formae binariae \((ac,ah,i)\ldots\varphi\). Haec forma certo non erit positiva (quum enim per hyp. \(a,b,c\) eadem signa non habeant, \(ab\) et \(ac\) simul positivi esse nequeunt); porro habebit numerum characteristicum \(-1\), quod synthetice ita demonstramus: Determinentur integri \(e,e'\) ita ut sit
\[
\begin{aligned}
e&\equiv0\pmod a\text{ et }\equiv\mathfrak B\pmod b;\\
ce'&\equiv\mathfrak A\pmod a\text{ et }\equiv h\mathfrak B\pmod b
\end{aligned}
\]
eritque \((e,e')\) valor expr. \(\surd-(ac,ah,i)\pmod{-ab}\). Nam secundum modulum \(a\) erit
\[
\begin{aligned}
ee&\equiv0\equiv-ac,\quad ee'\equiv0\equiv-ah\\
cce'e'&\equiv\mathfrak A\mathfrak A\equiv-bc\equiv-cci\quad\text{adeoque }e'e'\equiv-i
\end{aligned}
\]
secundum modulum \(b\) autem erit
\[
\begin{aligned}
ee&\equiv\mathfrak B\mathfrak B\equiv-ac,\quad cee'\equiv h\mathfrak B\mathfrak B\equiv-ach\quad\text{adeoque }ee'\equiv-ah\\
cce'e'&\equiv hh\mathfrak B\mathfrak B\equiv-achh\equiv-cci\quad\text{adeoque }e'e'\equiv-i
\end{aligned}
\]
eaedem vero tres congruentiae, quae secundum utrumque modulum \(a,b\) locum habent, etiam secundum modulum \(ab\) valebunt. Hinc per theoriam formarum ternariarum facile concluditur, \(\varphi\) repraesentabilem esse per formam \(\left(\begin{smallmatrix}-1,0,0\\1,0,0\end{smallmatrix}\right)\); sit itaque
\[
actt+2ahtu+iuu=-(\alpha t+\beta u)^2+2(\gamma t+\delta u)(\varepsilon t+\zeta u)
\]
\SourceLeaf{359}
eritque, multiplicando per \(c\),
\[
a(ct+hu)^2+buu=-c(\alpha t+\beta u)^2+2c(\gamma t+\delta u)(\varepsilon t+\zeta u)
\]
Hinc patet, si ipsis \(t,u\) tales valores determinati tribuantur, ut vel \(\gamma t+\delta u\), vel \(\varepsilon t+\zeta u\) fiat \(=0\), haberi solutionem aequationis \(\Omega\), cui igitur satisfiet tum per
\[
x=\delta c-\gamma h,\quad y=\gamma,\quad z=\alpha\delta-\beta\gamma
\]
tum per
\[
x=\zeta c-\varepsilon h,\quad y=\varepsilon,\quad z=\alpha\zeta-\beta\varepsilon
\]
simul manifestum est, neque illos valores neque hos simul \(=0\) fieri posse; si enim \(\delta c-\gamma h=0,\gamma=0\), fieret etiam \(\delta=0\) atque \(\varphi=-(\alpha t+\beta u)^2\), unde \(ab=0\) contra hyp., et perinde de alteris. --- In exemplo nostro invenimus formam \(\varphi\) hanc \((161,-63,24)\), valorem expr. \(\surd-\varphi\pmod{105}=(7,-51)\), atque repraesentationem formae \(\varphi\) per \(\left(\begin{smallmatrix}-1,0,0\\1,0,0\end{smallmatrix}\right)\) hanc,
\[
\varphi=-(13t-4u)^2+2(11t-4u)(15t-5u)
\]
hinc prodeunt solutiones \(x=7,y=11,z=-8\); \(x=20,y=15,z=-5\), sive dividendo per \(5\) et negligendo signum ipsius \(z\), \(x=4,y=3,z=1\).

Ex his duabus methodis aequationem \(\Omega\) solvendi posterior eo praestat, quod plerumque per numeros minores absolvitur; prior vero, quae etiam per varia artificia hic silentio praetereunda contrahi potest, elegantior videtur ea imprimis ratione, quod numeri \(a,b,c\) prorsus eodem modo tractantur, calculusque per horum permutationem quamcunque nihil mutatur. Hoc secus se habet in methodo secunda, ubi calculus maxime commodus plerumque provenit, si pro \(a\) accipitur minimus, pro \(c\) maximus trium numerorum datorum, uti in exemplo nostro fecimus.

\begin{center}\itshape De methodo per quam ill. Le Gendre theorema fundamentale tractavit.\end{center}
\Article{296}
Elegans theorema in artt. praecc. explicatum primo inventum est ab ill. Le Gendre, \emph{Hist. de l'Ac. de Paris} 1784 p. 507, atque demonstratione pulcra (a duabus nostris omnino diversa) munitum. Simul vero hic egregius geometra hoc loco operam dedit, demonstrationem propositionum, quae cum theoremate fundamentali
\SourceLeaf{360}
Sect. praec. conveniunt. inde derivare, quam ad hunc scopum non idoneam nobis videri iam supra declaravimus, art. 151. Hic itaque locus erit, hanc demonstrationem (per se valde elegantem) breviter exponendi iudiciique nostri rationes adiungendi. Praemittitur sequens observatio: \emph{Si numeri \(a,b,c\) omnes sunt \(\equiv1\pmod4\), aequatio \(axx+byy+czz=0\ldots(\Omega)\) solubilis esse nequit.} Facillime enim perspicitur, valorem ipsius \(axx+byy+czz\) necessario in hoc casu fieri vel \(\equiv1\), vel \(\equiv2\), vel \(\equiv3\pmod4\), nisi omnes \(x,y,z\) simul pares accipiantur; si itaque \(\Omega\) solubilis esset. hoc aliter fieri non posset quam per valores pares ipsorum \(x,y,z\), \emph{Q. E. A.}, quoniam valores quicunque aequationi \(\Omega\) satisfacientes etiamnum satisfaciunt, si per divisorem communem maximum dividuntur, unde necessario ad minimum unus impar prodire debet. Iam casus diversi theorematis demonstrandi ad sequentia momenta referuntur:

I. Designantibus \(p,q\) numeros primos formae \(4n+3\) (positivos inaequales), nequit simul esse \(p\,R\,q,q\,R\,p\). Si enim possibile esset, manifestum est statuendo \(1=a,-p=b,-q=c\), omnes conditiones ad resolubilitatem aequationis \(axx+byy+czz=0\) adimpletas esse (art. 294); eadem vero per observationem praec. resolutionem non admittit; quare suppositio consistere nequit. Hinc protinus sequitur propositio 7 art. 131.

II. Si \(p\) est numerus primus formae \(4n+1\), \(q\) numerus primus formae \(4n+3\), nequit simul esse \(q\,R\,p,p\,N\,q\). Alioquin enim foret \(-p\,R\,q\), atque aequatio \(xx+pyy-qzz=0\) resolubilis, quae per obs. praec. resolutionem respuit. Hinc derivantur casus 4 et 5 art. 131.

III. Si \(p,q\) sunt numeri primi formae \(4n+1\), nequit simul esse \(p\,R\,q,q\,N\,p\). Accipiatur alius numerus primus \(r\) formae \(4n+3\), qui sit residuum ipsius \(q\) et cuius non-residuum sit \(p\). Tunc erit per casus modo (II) demonstratos \(q\,R\,r,r\,N\,p\). Si itaque esset \(p\,R\,q,q\,N\,p\), foret \(qr\,R\,p,pr\,R\,q,pq\,N\,r\) et proin \(-pq\,R\,r\). Hinc aequatio \(pxx+qyy-rzz=0\) resolubilis esset contra obs. praec.; quare suppositio consistere nequit. Hinc sequuntur casus 1 et 2 art. 131.

Concinnius hic casus sequenti modo tractatur. Designet \(r\) numerum primum formae \(4n+3\), cuius non-residuum sit \(p\). Tunc erit etiam \(r\,N\,p\), adeoque (supponendo \(p\,R\,q,q\,N\,p\)) \(qr\,R\,p\). porro \(-p\,R\,q,-p\,R\,r\) et proin etiam \(-p\,R\,qr\); quare aequatio \(xx+pyy-qrzz=0\) resolubilis esset contra obs. praec. Hinc etc.

IV. Si \(p\) est numerus primus formae \(4n+1\), \(q\) primus formae \(4n+3\),
\SourceLeaf{361}
nequit simul esse \(p\,R\,q,q\,N\,p\). Accipiatur numerus primus auxiliaris \(r\) formae \(4n+1\) qui sit non-residuum utriusque \(p,q\). Tunc erit (per II) \(q\,N\,r\) et (per III) \(p\,N\,r\); hinc \(pq\,R\,r\); si itaque esset \(p\,R\,q,q\,N\,p\), haberetur etiam \(pr\,N\,q,-pr\,R\,q,qr\,R\,p\); quare aequatio \(pxx-qyy+rzz=0\) resolubilis esset, \emph{Q. E. A.} --- Hinc derivantur casus 3 et 6 art. 131.

V. Designantibus \(p,q\) numeros primos formae \(4n+3\), nequit simul esse \(p\,N\,q,q\,N\,p\). Supponendo enim fieri posse, et accipiendo numerum primum auxiliarem \(r\) formae \(4n+1\), qui sit non-residuum utriusque \(p,q\) erit \(qr\,R\,p,pr\,R\,q\); porro (per II) \(p\,N\,r,q\,N\,r\), unde \(pq\,R\,r\) et \(-pq\,R\,r\); hinc aequatio \(-pxx-qyy+rzz=0\) possibilis, contra obs. praec. Hinc deducitur casus 8 art. 131.

\Article{297}
Demonstrationem praec. proprius contemplando quisque facile intelliget, casus I et II ita absolutos esse ut nihil obiici possit. At demonstrationes casuum reliquorum innituntur existentiae numerorum auxiliarium, qua nondum demonstrata methodus manifesto omnem vim perdit. Quae suppositiones, etsi tam speciosae sint, ut minus attendenti demonstratione ne opus quidem esse videri possit, atque certe theorema demonstrandum ad maximum \emph{probabilitatis} gradum evehant, tamen si rigor geometricus desideretur, neutiquam gratuito sunt admittendae. Quod quidem attinet ad suppositionem in IV et V, exstare numerum primum \(r\) formae \(4n+1\), qui duorum aliorum primorum datorum \(p,q\) non-residuum sit, e Sect. IV facile concluditur, omnes numeros ipso \(4pq\) minores ad ipsumque primos (quorum multitudo est \(2(p-1)(q-1)\)) in quatuor classes aequaliter distribui, quarum una contineat non-residua utriusque \(p,q\), tres reliquae residua ipsius \(p\) non-residua ipsius \(q\), non-residua ipsius \(p\) residua ipsius \(q\), residua utriusque \(p,q\); et in singulis classibus semissem fore numeros formae \(4n+1\), semissem formae \(4n+3\). Habebuntur itaque inter illos \(\frac14(p-1)(q-1)\) non-residua utriusque \(p,q\) formae \(4n+1\), qui sint \(g.g',g''\) etc.; numeri \(\frac74(p-1)(q-1)\) reliqui sint \(h,h',h''\) etc. Manifesto omnes numeri in formis \(4pqt+g,4pqt+g',4pqt+g''\) etc. \((G)\) contenti quoque erunt non-residua ipsorum \(p,q\) formae \(4n+1\). Iam patet, ad suppositionem stabiliendam demonstrari tantummodo debere, sub formis \((G)\) certo contineri \emph{numeros primos}, quod sane iam per se valde plausibile videtur, quum hae formae una cum his \(4pqt+h,4pqt+h'\) etc.
